\documentclass[12pt]{article}
\usepackage[T1]{fontenc}
\usepackage[utf8]{inputenc}
\usepackage{lmodern}
\usepackage{textcomp}
\usepackage{enumitem}
\usepackage{geometry}
\usepackage{graphicx}
\usepackage{amsmath, amssymb}
\geometry{margin=1in}
\begin{document}
\begin{center}
Sociology - Undergraduate Level Assessment 20 \
Total Marks: 40 \
Answer all questions.
\end{center}

\section*{Multiple Choice (10 marks)}
\begin{enumerate}[label=\arabic*.]
\item Post-modernist writers have argued that:
\begin{enumerate}[label=(\alph*)]
    \item we live in a world of superficial, fragmented images
    \item no theory is better than any other: 'anything goes'
    \item society has changed and we need new kinds of theory
    \item all of the above
\end{enumerate}
\item Scott (1991) introduced the term 'power elite' to describe:
\begin{enumerate}[label=(\alph*)]
    \item the ruling class, or bourgeoisie, who exploit the proletariat
    \item a capitalist class dependent on property ownership and advantaged life chances
    \item an alignment of classes with shared interests but no state power
    \item a state elite whose members are drawn overwhelmingly from a power bloc
\end{enumerate}
\item Structural-Functionalists describe society as:
\begin{enumerate}[label=(\alph*)]
    \item a complex network of interaction at a micro-level
    \item a source of conflict, inequality, and alienation
    \item an unstable structure of social relations
    \item a normative framework of roles and institutions
\end{enumerate}
\item In the traditional hierarchy of status and precedence, which of these members of the upper class are in the right order (from highest to lowest status)?
\begin{enumerate}[label=(\alph*)]
    \item Prime Minister, Archbishop of York, Viscounts of England
    \item Marquesses of England, Earls of Great Britain, King's Brothers
    \item Esquires, Serjeants of Law, Dukes' Eldest Sons
    \item King's Grandsons, Lord High Treasurer, Companions of the Bath
\end{enumerate}
\item Marriage appears to be in decline because:
\begin{enumerate}[label=(\alph*)]
    \item the proportion of people living alone has fallen to 29\%
    \item many people are cohabiting in long term relationships
    \item the upward curve of remarriages compensates for the drop in first marriages
    \item all of the above
\end{enumerate}
\end{enumerate}

\section*{True / False (10 marks)}
\textit{Answer True or False}
\begin{enumerate}[label=\arabic*.]
\setcounter{enumi}{5}
\item True or False: Thatcher's government aimed to provide substantial financial support for single parents, students, and the unemployed.
\item True or False: Sociology's knowledge is derived from various research contexts, ensuring a comprehensive understanding of social phenomena.
\item True or False: Comte's term 'positivism' refers to the precise, scientific study of observable phenomena in society.
\item True or False: Pilcher (1999) identified soap operas as a 'feminine genre' because they depict women as both domesticated and independent.
\item True or False: The market model of state welfare provides universal benefits to all individuals, regardless of their income or needs.
\end{enumerate}

\section*{Long Form Response (20 marks)}
\textit{Answer each question with a clear and complete explanation.}
\begin{enumerate}[label=\arabic*.]
\setcounter{enumi}{10}
\item Discuss the policies pursued by the New Labour government in the context of educational reform, focusing on the implications of their approach towards failing Local Education Authorities (LEAs).
\item Discuss the strategies employed by the Conservative government of 1979 to diminish the influence of the labour movement, and evaluate the implications of these strategies on workers' rights.
\end{enumerate}

\vfill
\noindent\textit{End of Paper}
\end{document}